% Định nghĩa lệnh Lý thuyết — dùng khi biên dịch lt.tex bằng TeX.
% App web đọc lt.tex trực tiếp (bỏ \input file này).
% Trong lt.tex, từ thư mục bài:
%   % Định nghĩa lệnh Lý thuyết — dùng khi biên dịch lt.tex bằng TeX.
% App web đọc lt.tex trực tiếp (bỏ \input file này).
% Trong lt.tex, từ thư mục bài:
%   % Định nghĩa lệnh Lý thuyết — dùng khi biên dịch lt.tex bằng TeX.
% App web đọc lt.tex trực tiếp (bỏ \input file này).
% Trong lt.tex, từ thư mục bài:
%   \input{../../../../_lenh/lythuyet.tex}
% từ gốc repo:
%   \input{ngan-hang/_lenh/lythuyet.tex}
\makeatletter
\providecommand{\indam}[1]{\textbf{#1}}
\providecommand{\chuy}[1]{\par\smallskip\noindent\textit{#1}\par}
\providecommand{\luuy}[1]{\par\smallskip\noindent\textbf{Lưu ý.} #1\par}
\providecommand{\ghichu}[1]{\par\smallskip\noindent\textbf{Ghi chú.} #1\par}
\providecommand{\vidu}[1]{\par\smallskip\noindent\textbf{Ví dụ.} #1\par}
\providecommand{\Blythuyet}{\subsection*{Lý thuyết}}
% Hai cột: kiến thức | hình
\providecommand{\haicotchay}[2]{%
  \par\medskip\noindent
  \begin{minipage}[t]{0.52\linewidth}#1\end{minipage}\hfill
  \begin{minipage}[t]{0.46\linewidth}#2\end{minipage}%
  \par\medskip
}
\providecommand{\kienthuccot}[2]{\haicotchay{#1}{#2}}
% Dự phòng nếu chưa nạp graphicx / caption (không ghi đè bản thật)
\providecommand{\resizebox}[3]{#3}
\providecommand{\captionof}[2]{\par\smallskip\noindent\textit{#2}\par}
\newenvironment{hoatdong}[1]{%
  \par\medskip\noindent\textbf{Hoạt động.} #1\par\smallskip
}{\par\medskip}
\newenvironment{emcobiet}{%
  \par\medskip\noindent\textbf{Em có biết}\par\smallskip
}{\par\medskip}
\newenvironment{emdahoc}{%
  \par\medskip\noindent\textbf{Em đã học}\par\smallskip
}{\par\medskip}
\newenvironment{emcothe}{%
  \par\medskip\noindent\textbf{Em có thể}\par\smallskip
}{\par\medskip}
\newenvironment{kienthuc}{%
  \par\medskip\noindent\textbf{Kiến thức cốt lõi}\par\smallskip
}{\par\medskip}
\newenvironment{traloi}{%
  \par\smallskip\noindent\textit{Trả lời / ghi chú của em}\par
  \vspace{1.2em}%
}{\par\medskip}
\makeatother

% từ gốc repo:
%   % Định nghĩa lệnh Lý thuyết — dùng khi biên dịch lt.tex bằng TeX.
% App web đọc lt.tex trực tiếp (bỏ \input file này).
% Trong lt.tex, từ thư mục bài:
%   \input{../../../../_lenh/lythuyet.tex}
% từ gốc repo:
%   \input{ngan-hang/_lenh/lythuyet.tex}
\makeatletter
\providecommand{\indam}[1]{\textbf{#1}}
\providecommand{\chuy}[1]{\par\smallskip\noindent\textit{#1}\par}
\providecommand{\luuy}[1]{\par\smallskip\noindent\textbf{Lưu ý.} #1\par}
\providecommand{\ghichu}[1]{\par\smallskip\noindent\textbf{Ghi chú.} #1\par}
\providecommand{\vidu}[1]{\par\smallskip\noindent\textbf{Ví dụ.} #1\par}
\providecommand{\Blythuyet}{\subsection*{Lý thuyết}}
% Hai cột: kiến thức | hình
\providecommand{\haicotchay}[2]{%
  \par\medskip\noindent
  \begin{minipage}[t]{0.52\linewidth}#1\end{minipage}\hfill
  \begin{minipage}[t]{0.46\linewidth}#2\end{minipage}%
  \par\medskip
}
\providecommand{\kienthuccot}[2]{\haicotchay{#1}{#2}}
% Dự phòng nếu chưa nạp graphicx / caption (không ghi đè bản thật)
\providecommand{\resizebox}[3]{#3}
\providecommand{\captionof}[2]{\par\smallskip\noindent\textit{#2}\par}
\newenvironment{hoatdong}[1]{%
  \par\medskip\noindent\textbf{Hoạt động.} #1\par\smallskip
}{\par\medskip}
\newenvironment{emcobiet}{%
  \par\medskip\noindent\textbf{Em có biết}\par\smallskip
}{\par\medskip}
\newenvironment{emdahoc}{%
  \par\medskip\noindent\textbf{Em đã học}\par\smallskip
}{\par\medskip}
\newenvironment{emcothe}{%
  \par\medskip\noindent\textbf{Em có thể}\par\smallskip
}{\par\medskip}
\newenvironment{kienthuc}{%
  \par\medskip\noindent\textbf{Kiến thức cốt lõi}\par\smallskip
}{\par\medskip}
\newenvironment{traloi}{%
  \par\smallskip\noindent\textit{Trả lời / ghi chú của em}\par
  \vspace{1.2em}%
}{\par\medskip}
\makeatother

\makeatletter
\providecommand{\indam}[1]{\textbf{#1}}
\providecommand{\chuy}[1]{\par\smallskip\noindent\textit{#1}\par}
\providecommand{\luuy}[1]{\par\smallskip\noindent\textbf{Lưu ý.} #1\par}
\providecommand{\ghichu}[1]{\par\smallskip\noindent\textbf{Ghi chú.} #1\par}
\providecommand{\vidu}[1]{\par\smallskip\noindent\textbf{Ví dụ.} #1\par}
\providecommand{\Blythuyet}{\subsection*{Lý thuyết}}
% Hai cột: kiến thức | hình
\providecommand{\haicotchay}[2]{%
  \par\medskip\noindent
  \begin{minipage}[t]{0.52\linewidth}#1\end{minipage}\hfill
  \begin{minipage}[t]{0.46\linewidth}#2\end{minipage}%
  \par\medskip
}
\providecommand{\kienthuccot}[2]{\haicotchay{#1}{#2}}
% Dự phòng nếu chưa nạp graphicx / caption (không ghi đè bản thật)
\providecommand{\resizebox}[3]{#3}
\providecommand{\captionof}[2]{\par\smallskip\noindent\textit{#2}\par}
\newenvironment{hoatdong}[1]{%
  \par\medskip\noindent\textbf{Hoạt động.} #1\par\smallskip
}{\par\medskip}
\newenvironment{emcobiet}{%
  \par\medskip\noindent\textbf{Em có biết}\par\smallskip
}{\par\medskip}
\newenvironment{emdahoc}{%
  \par\medskip\noindent\textbf{Em đã học}\par\smallskip
}{\par\medskip}
\newenvironment{emcothe}{%
  \par\medskip\noindent\textbf{Em có thể}\par\smallskip
}{\par\medskip}
\newenvironment{kienthuc}{%
  \par\medskip\noindent\textbf{Kiến thức cốt lõi}\par\smallskip
}{\par\medskip}
\newenvironment{traloi}{%
  \par\smallskip\noindent\textit{Trả lời / ghi chú của em}\par
  \vspace{1.2em}%
}{\par\medskip}
\makeatother

% từ gốc repo:
%   % Định nghĩa lệnh Lý thuyết — dùng khi biên dịch lt.tex bằng TeX.
% App web đọc lt.tex trực tiếp (bỏ \input file này).
% Trong lt.tex, từ thư mục bài:
%   % Định nghĩa lệnh Lý thuyết — dùng khi biên dịch lt.tex bằng TeX.
% App web đọc lt.tex trực tiếp (bỏ \input file này).
% Trong lt.tex, từ thư mục bài:
%   \input{../../../../_lenh/lythuyet.tex}
% từ gốc repo:
%   \input{ngan-hang/_lenh/lythuyet.tex}
\makeatletter
\providecommand{\indam}[1]{\textbf{#1}}
\providecommand{\chuy}[1]{\par\smallskip\noindent\textit{#1}\par}
\providecommand{\luuy}[1]{\par\smallskip\noindent\textbf{Lưu ý.} #1\par}
\providecommand{\ghichu}[1]{\par\smallskip\noindent\textbf{Ghi chú.} #1\par}
\providecommand{\vidu}[1]{\par\smallskip\noindent\textbf{Ví dụ.} #1\par}
\providecommand{\Blythuyet}{\subsection*{Lý thuyết}}
% Hai cột: kiến thức | hình
\providecommand{\haicotchay}[2]{%
  \par\medskip\noindent
  \begin{minipage}[t]{0.52\linewidth}#1\end{minipage}\hfill
  \begin{minipage}[t]{0.46\linewidth}#2\end{minipage}%
  \par\medskip
}
\providecommand{\kienthuccot}[2]{\haicotchay{#1}{#2}}
% Dự phòng nếu chưa nạp graphicx / caption (không ghi đè bản thật)
\providecommand{\resizebox}[3]{#3}
\providecommand{\captionof}[2]{\par\smallskip\noindent\textit{#2}\par}
\newenvironment{hoatdong}[1]{%
  \par\medskip\noindent\textbf{Hoạt động.} #1\par\smallskip
}{\par\medskip}
\newenvironment{emcobiet}{%
  \par\medskip\noindent\textbf{Em có biết}\par\smallskip
}{\par\medskip}
\newenvironment{emdahoc}{%
  \par\medskip\noindent\textbf{Em đã học}\par\smallskip
}{\par\medskip}
\newenvironment{emcothe}{%
  \par\medskip\noindent\textbf{Em có thể}\par\smallskip
}{\par\medskip}
\newenvironment{kienthuc}{%
  \par\medskip\noindent\textbf{Kiến thức cốt lõi}\par\smallskip
}{\par\medskip}
\newenvironment{traloi}{%
  \par\smallskip\noindent\textit{Trả lời / ghi chú của em}\par
  \vspace{1.2em}%
}{\par\medskip}
\makeatother

% từ gốc repo:
%   % Định nghĩa lệnh Lý thuyết — dùng khi biên dịch lt.tex bằng TeX.
% App web đọc lt.tex trực tiếp (bỏ \input file này).
% Trong lt.tex, từ thư mục bài:
%   \input{../../../../_lenh/lythuyet.tex}
% từ gốc repo:
%   \input{ngan-hang/_lenh/lythuyet.tex}
\makeatletter
\providecommand{\indam}[1]{\textbf{#1}}
\providecommand{\chuy}[1]{\par\smallskip\noindent\textit{#1}\par}
\providecommand{\luuy}[1]{\par\smallskip\noindent\textbf{Lưu ý.} #1\par}
\providecommand{\ghichu}[1]{\par\smallskip\noindent\textbf{Ghi chú.} #1\par}
\providecommand{\vidu}[1]{\par\smallskip\noindent\textbf{Ví dụ.} #1\par}
\providecommand{\Blythuyet}{\subsection*{Lý thuyết}}
% Hai cột: kiến thức | hình
\providecommand{\haicotchay}[2]{%
  \par\medskip\noindent
  \begin{minipage}[t]{0.52\linewidth}#1\end{minipage}\hfill
  \begin{minipage}[t]{0.46\linewidth}#2\end{minipage}%
  \par\medskip
}
\providecommand{\kienthuccot}[2]{\haicotchay{#1}{#2}}
% Dự phòng nếu chưa nạp graphicx / caption (không ghi đè bản thật)
\providecommand{\resizebox}[3]{#3}
\providecommand{\captionof}[2]{\par\smallskip\noindent\textit{#2}\par}
\newenvironment{hoatdong}[1]{%
  \par\medskip\noindent\textbf{Hoạt động.} #1\par\smallskip
}{\par\medskip}
\newenvironment{emcobiet}{%
  \par\medskip\noindent\textbf{Em có biết}\par\smallskip
}{\par\medskip}
\newenvironment{emdahoc}{%
  \par\medskip\noindent\textbf{Em đã học}\par\smallskip
}{\par\medskip}
\newenvironment{emcothe}{%
  \par\medskip\noindent\textbf{Em có thể}\par\smallskip
}{\par\medskip}
\newenvironment{kienthuc}{%
  \par\medskip\noindent\textbf{Kiến thức cốt lõi}\par\smallskip
}{\par\medskip}
\newenvironment{traloi}{%
  \par\smallskip\noindent\textit{Trả lời / ghi chú của em}\par
  \vspace{1.2em}%
}{\par\medskip}
\makeatother

\makeatletter
\providecommand{\indam}[1]{\textbf{#1}}
\providecommand{\chuy}[1]{\par\smallskip\noindent\textit{#1}\par}
\providecommand{\luuy}[1]{\par\smallskip\noindent\textbf{Lưu ý.} #1\par}
\providecommand{\ghichu}[1]{\par\smallskip\noindent\textbf{Ghi chú.} #1\par}
\providecommand{\vidu}[1]{\par\smallskip\noindent\textbf{Ví dụ.} #1\par}
\providecommand{\Blythuyet}{\subsection*{Lý thuyết}}
% Hai cột: kiến thức | hình
\providecommand{\haicotchay}[2]{%
  \par\medskip\noindent
  \begin{minipage}[t]{0.52\linewidth}#1\end{minipage}\hfill
  \begin{minipage}[t]{0.46\linewidth}#2\end{minipage}%
  \par\medskip
}
\providecommand{\kienthuccot}[2]{\haicotchay{#1}{#2}}
% Dự phòng nếu chưa nạp graphicx / caption (không ghi đè bản thật)
\providecommand{\resizebox}[3]{#3}
\providecommand{\captionof}[2]{\par\smallskip\noindent\textit{#2}\par}
\newenvironment{hoatdong}[1]{%
  \par\medskip\noindent\textbf{Hoạt động.} #1\par\smallskip
}{\par\medskip}
\newenvironment{emcobiet}{%
  \par\medskip\noindent\textbf{Em có biết}\par\smallskip
}{\par\medskip}
\newenvironment{emdahoc}{%
  \par\medskip\noindent\textbf{Em đã học}\par\smallskip
}{\par\medskip}
\newenvironment{emcothe}{%
  \par\medskip\noindent\textbf{Em có thể}\par\smallskip
}{\par\medskip}
\newenvironment{kienthuc}{%
  \par\medskip\noindent\textbf{Kiến thức cốt lõi}\par\smallskip
}{\par\medskip}
\newenvironment{traloi}{%
  \par\smallskip\noindent\textit{Trả lời / ghi chú của em}\par
  \vspace{1.2em}%
}{\par\medskip}
\makeatother

\makeatletter
\providecommand{\indam}[1]{\textbf{#1}}
\providecommand{\chuy}[1]{\par\smallskip\noindent\textit{#1}\par}
\providecommand{\luuy}[1]{\par\smallskip\noindent\textbf{Lưu ý.} #1\par}
\providecommand{\ghichu}[1]{\par\smallskip\noindent\textbf{Ghi chú.} #1\par}
\providecommand{\vidu}[1]{\par\smallskip\noindent\textbf{Ví dụ.} #1\par}
\providecommand{\Blythuyet}{\subsection*{Lý thuyết}}
% Hai cột: kiến thức | hình
\providecommand{\haicotchay}[2]{%
  \par\medskip\noindent
  \begin{minipage}[t]{0.52\linewidth}#1\end{minipage}\hfill
  \begin{minipage}[t]{0.46\linewidth}#2\end{minipage}%
  \par\medskip
}
\providecommand{\kienthuccot}[2]{\haicotchay{#1}{#2}}
% Dự phòng nếu chưa nạp graphicx / caption (không ghi đè bản thật)
\providecommand{\resizebox}[3]{#3}
\providecommand{\captionof}[2]{\par\smallskip\noindent\textit{#2}\par}
\newenvironment{hoatdong}[1]{%
  \par\medskip\noindent\textbf{Hoạt động.} #1\par\smallskip
}{\par\medskip}
\newenvironment{emcobiet}{%
  \par\medskip\noindent\textbf{Em có biết}\par\smallskip
}{\par\medskip}
\newenvironment{emdahoc}{%
  \par\medskip\noindent\textbf{Em đã học}\par\smallskip
}{\par\medskip}
\newenvironment{emcothe}{%
  \par\medskip\noindent\textbf{Em có thể}\par\smallskip
}{\par\medskip}
\newenvironment{kienthuc}{%
  \par\medskip\noindent\textbf{Kiến thức cốt lõi}\par\smallskip
}{\par\medskip}
\newenvironment{traloi}{%
  \par\smallskip\noindent\textit{Trả lời / ghi chú của em}\par
  \vspace{1.2em}%
}{\par\medskip}
\makeatother


\subsection{Lý thuyết}

\subsubsection{An toàn khi sử dụng thiết bị thí nghiệm}
\paragraph{Sử dụng các thiết bị điện}
\begin{itemize}
	\item Các thiết bị thí nghiệm Vật lí phổ thông sử dụng nguồn điện đa dạng như máy biến áp (máy biến thế), bộ chuyển đổi điện áp AC-DC để cung cấp dòng điện một chiều và xoay chiều ở các mức điện áp khác nhau.
	\item Một số kí hiệu và nhãn thông số thông dụng trên thiết bị thí nghiệm cần phải quan sát để sử dụng đúng chức năng, đúng yêu cầu kĩ thuật:
	\begin{itemize}
		\item DC hoặc dấu "$-$" biểu diễn dòng điện một chiều.
		\item AC hoặc dấu "$\sim$" biểu diễn dòng điện xoay chiều.
		\item Input (I) chỉ đầu vào, Output (O) chỉ đầu ra của nguồn điện hoặc thiết bị.
		\item Cực dương được kí hiệu là "$+$" hoặc sơn màu đỏ; cực âm được kí hiệu là "$-$" hoặc sơn màu xanh (hoặc màu đen).
		\item Biểu tượng hai mũi tên hướng lên biểu thị dụng cụ cần đặt đứng thẳng.
		\item Biểu tượng chiếc ly bị nứt biểu thị dụng cụ dễ vỡ, cần nhẹ tay.
		\item Biểu tượng mặt trời tỏa sáng biểu thị thiết bị cần tránh ánh nắng chiếu trực tiếp.
		\item Biểu tượng thùng rác gạch chéo biểu thị thiết bị không được phép bỏ vào thùng rác sinh hoạt thông thường.
	\end{itemize}
	\item Để đảm bảo an toàn khi sử dụng các thiết bị điện, người dùng cần tuân thủ nghiêm ngặt các hướng dẫn:
	\begin{itemize}
		\item Chỉ cắm phích cắm của thiết bị điện vào ổ cắm khi hiệu điện thế của nguồn tương ứng với hiệu điện thế định mức của dụng cụ.
		\item Luôn tắt công tắc nguồn của thiết bị điện trước khi thực hiện cắm hoặc tháo phích cắm điện.
		\item Phải bố trí dây dẫn điện gọn gàng, không bị chồng chéo hay vướng mắc khi đi lại trong phòng thực hành.
	\end{itemize}
\end{itemize}

\paragraph{Sử dụng các thiết bị nhiệt và thuỷ tinh}
\begin{itemize}
	\item Các thiết bị đun nóng (đèn cồn, lò nung, bếp hồng ngoại) có thể gây bỏng nghiêm trọng cho người sử dụng, làm nứt vỡ các bộ phận làm bằng thuỷ tinh hoặc gây hỏa hoạn.
	\item Tuyệt đối không tiếp xúc trực tiếp với các vật và các thiết bị thí nghiệm đang có nhiệt độ cao khi không có các dụng cụ bảo hộ chịu nhiệt chuyên dụng.
	\item Không được để nước cũng như các dung dịch dẫn điện hoặc dung dịch dễ cháy gần các thiết bị điện đang hoạt động.
	\item Giữ khoảng cách an toàn khi tiến hành các thí nghiệm nung nóng các vật để tránh hỏa hoạn hoặc bỏng do vật bắn ra.
\end{itemize}

\paragraph{Sử dụng các thiết bị quang học}
\begin{itemize}
	\item Các thiết bị quang học như kính lăng, thấu kính, gương phẳng rất tinh xảo, dễ bị mốc, xước, nứt vỡ hoặc dính bụi bẩn làm sai lệch đường truyền của tia sáng và ảnh hưởng trực tiếp đến độ chính xác của kết quả thí nghiệm.
	\item Tránh dùng tay tiếp xúc trực tiếp vào bề mặt của thấu kính, gương, lăng kính; cần bảo quản thiết bị trong hộp chuyên dụng sau khi dùng xong.
	\item Tia laser có mật độ năng lượng rất cao, có khả năng làm tổn thương nghiêm trọng hoặc gây mù võng mạc mắt. Vì vậy, tuyệt đối không nhìn trực tiếp vào chùm tia laser hoặc hướng chùm tia laser vào mắt người khác.
\end{itemize}

\subsubsection{Nguy cơ mất an toàn trong sử dụng thiết bị thí nghiệm Vật lí}

\paragraph{Nguy cơ gây nguy hiểm cho người sử dụng}
\begin{itemize}
\item Thao tác cắm phích điện vào ổ cắm lỏng lẻo, dùng tay ướt để cắm điện hoặc cầm vào dây điện để giật mạnh phích cắm thay vì cầm vào vỏ nhựa cách điện của phích cắm là những nguyên nhân hàng đầu gây điện giật.
\item Đun nước trên đèn cồn không cẩn thận, để đèn cồn bị đổ hoặc để các vật liệu dễ bắt lửa gần ngọn lửa đèn cồn sẽ dẫn đến nguy cơ cháy nổ nghiêm trọng.
\item Sử dụng các nguồn phóng xạ hoặc tia laser không đúng quy định che chắn cơ thể, không đeo kính bảo hộ chuyên dụng có thể dẫn đến nhiễm xạ tế bào hoặc hỏng thị lực.
\end{itemize}

\paragraph{Nguy cơ hỏng thiết bị đo điện}
\begin{itemize}
\item Sử dụng ampe kế hoặc vôn kế để đo dòng điện hoặc điện áp vượt quá giới hạn đo (thang đo) của thiết bị sẽ làm kim đo bị kịch khung hoặc làm cuộn dây bên trong bị quá tải, gây chập cháy thiết bị đo.
\item Chọn sai chức năng đo (ví dụ: cắm nhầm dây đo vào thang đo dòng điện nhưng lại đo hiệu điện thế) sẽ tạo ra một mạch ngắn mạch cực mạnh làm hỏng toàn bộ mạch điện và cháy đồng hồ đo đa năng.
\end{itemize}

\paragraph{Nguy cơ cháy nổ trong phòng thực hành}
\begin{itemize}
\item Để các dung dịch dễ cháy bốc hơi (như cồn, ether, xăng) gần nguồn nhiệt hoặc thí nghiệm mạch điện có nguy cơ phát tia lửa điện là nguồn cơn gây ra các vụ nổ và hỏa hoạn.
\item Khi phát hiện phòng thực hành xảy ra đám cháy, cần ngay lập tức ngắt cầu dao tổng của hệ thống điện, báo động cho mọi người xung quanh, di chuyển các chất dễ cháy ra khỏi vùng nguy hiểm và sử dụng các biện pháp dập lửa phù hợp (bình cứu hỏa CO2, cát khô hoặc bình bột). Tuyệt đối không dùng nước để dập các đám cháy điện hoặc đám cháy kim loại hoạt động mạnh như magnesium, natri.
\end{itemize}

\subsubsection{Quy tắc an toàn trong phòng thực hành}

\paragraph{Biển báo cảnh báo và biển chỉ dẫn trong phòng thí nghiệm}
\begin{itemize}
\item Các biển báo an toàn trong phòng thực hành Vật lí được chia làm 3 loại chính:
\begin{itemize}
\item Biển cảnh báo cấm (hình tròn viền đỏ, nền trắng, hình vẽ màu đen): có ý nghĩa cấm thực hiện các hành động có thể gây nguy hiểm như cấm lửa, cấm ăn uống, cấm nguồn nước uống trực tiếp...
\item Biển cảnh báo nguy hiểm (hình tam giác đều viền đen, nền vàng, hình vẽ màu đen): cảnh báo các mối nguy hiểm tiềm tàng như nơi nguy hiểm về điện, từ trường cao, nơi có chất phóng xạ, cảnh báo tia laser, cảnh báo chất ăn mòn, cảnh báo nhiệt độ cao, cảnh báo vật sắc nhọn...
\item Biển bắt buộc thực hiện hoặc chỉ dẫn (hình chữ nhật hoặc hình tròn nền xanh dương/xanh lá, hình vẽ màu trắng): chỉ dẫn lối thoát hiểm, vị trí bình cứu hỏa hoặc yêu cầu bắt buộc mang đồ bảo hộ như kính bảo hộ mắt, mặt nạ phòng độc, găng tay chống hóa chất, áo choàng bảo hộ...
\end{itemize}
\end{itemize}

\paragraph{Quy tắc xử lí sự cố mất an toàn thực tế}
\begin{itemize}
\item Khi không may bị hoá chất ăn da hoặc axit bám vào tay, bước đầu tiên và quan trọng nhất là phải lập tức xả tay dưới vòi nước sạch chảy liên tục trong ít nhất 10–15 phút để làm loãng hóa chất, sau đó báo cáo ngay với giáo viên hướng dẫn.
\item Khi xảy ra sự cố vỡ nhiệt kế thuỷ ngân, tuyệt đối không được dùng chổi để quét dọn mạnh vì sẽ làm thủy ngân bị phân tán thành nhiều hạt siêu nhỏ, tăng tốc độ bay hơi gây nhiễm độc hô hấp. Cách xử lí đúng là: sơ tán mọi người ra khỏi khu vực, đeo khẩu trang và găng tay cao su, rắc một lớp bột lưu huỳnh lên những hạt thủy ngân thu được để tạo thành hợp chất đồng sunfua không bay hơi (\\(HgS\\)), hoặc dùng tấm giấy mỏng gom các giọt thủy ngân lại rồi bỏ vào lọ kín nút chặt.
\item Để đảm bảo an toàn tuyệt đối khi làm việc với các nguồn phóng xạ, con người phải triệt để áp dụng 3 quy tắc: giảm tối đa thời gian tiếp xúc, tăng tối đa khoảng cách từ cơ thể đến nguồn phóng xạ, và bắt buộc phải che chắn cơ thể bằng các tấm chắn chì hoặc mặc áo bảo hộ bằng chì chuyên dụng.
\end{itemize}

\begin{kienthuc}
Các ý kiến thức cốt lõi của bài:
\begin{itemize}
\item Hiểu rõ và phân biệt được ý nghĩa các kí hiệu điện áp, các nhãn thông số kĩ thuật trên các thiết bị đo điện, nguồn điện AC-DC thông dụng.
\item Nắm vững các nguy cơ gây cháy nổ, giật điện, hỏng hóc thiết bị đo và biết cách phòng ngừa các tai nạn này trong quá trình làm thực hành Vật lí.
\item Phân loại được 3 nhóm biển báo cảnh báo an toàn cơ bản và nêu được công dụng của từng trang thiết bị bảo hộ cá nhân (kính, găng tay, áo chì).
\item Thực hiện thành thạo và chuẩn xác các thao tác xử lí sự cố thực tế thường gặp như sơ cứu khi dính hóa chất ăn da, xử lí nhiệt kế thủy ngân bị vỡ hay dập tắt đám cháy đúng phương pháp.
\end{itemize}
\end{kienthuc}

\begin{emdahoc}
\begin{itemize}
\item Cần quan sát, đọc kĩ các kí hiệu và nhãn thông số kĩ thuật trên các thiết bị thí nghiệm trước khi kết nối vào nguồn điện để bảo vệ an toàn cho cả người sử dụng và thiết bị.
\item Phải tuân thủ nghiêm ngặt các quy tắc phòng cháy chữa cháy, an toàn điện, an toàn nhiệt, quang học và an toàn khi làm việc với các chất độc hại hoặc nguồn phóng xạ trong phòng thí nghiệm.
\item Luôn thực hiện thí nghiệm dưới sự cho phép và giám sát trực tiếp của giáo viên hướng dẫn, báo cáo ngay lập tức khi phát hiện sự cố mất an toàn.
\end{itemize}
\end{emdahoc}

\begin{emcothe}
\begin{itemize}
\item Giải thích được nguyên nhân gây nguy hiểm về điện, hỏa hoạn từ đó đề xuất được các giải pháp, nội quy thiết thực để xây dựng môi trường phòng thực hành Vật lí an toàn.
\item Thiết kế và thuyết trình được bảng hướng dẫn các quy tắc an toàn hoặc poster tuyên truyền phòng chống nguy hiểm phóng xạ, tia laser trong đời sống và sản xuất.
\item Chủ động xử lí thông minh, an toàn các tình huống khẩn cấp xảy ra trong gia đình hoặc trường học liên quan đến chập cháy điện hoặc rò rỉ hóa chất độc hại.
\end{itemize}
\end{emcothe}
